% concatenated from main.tex, Case_A7_B5.tex, Case_A7_B6_A1_3.tex, Case_A7_B6_A1_2.tex, Case_A7_B6_A1_1.tex, appendix.tex
\documentclass[a4paper,fleqn]{cas-sc}

\usepackage[numbers,sort&compress]{natbib}

\usepackage{amsmath,amssymb,amsfonts,amsthm,mathtools}
\usepackage{booktabs}
\usepackage{tikz}
\usepackage{tabularx}
\usepackage{xurl}

\usepackage{enumitem}
\usepackage{array}
%\usepackage{lineno}
%\linenumbers

\usepackage{hyperref}
\usepackage{microtype}
\emergencystretch=3em


\newtheorem{theorem}{Theorem}[section]
\newtheorem{lemma}[theorem]{Lemma}
\newtheorem{proposition}[theorem]{Proposition}
\newtheorem{corollary}[theorem]{Corollary}
\newtheorem{claim}[theorem]{Claim}
\newtheorem{conjecture}[theorem]{Conjecture}

\theoremstyle{definition}
\newtheorem{definition}[theorem]{Definition}
\newtheorem{example}[theorem]{Example}
\newtheorem{remark}[theorem]{Remark}
\newtheorem{computationallemma}{Computational Lemma}

\DeclareMathOperator{\dist}{dist}
\DeclareMathOperator{\diam}{diam}
\DeclareMathOperator{\OPT}{OPT}
\newcommand{\sas}[1]{\textcolor{blue}{SS: #1}}
\newcommand{\arp}[1]{\textcolor{red}{ArP: #1}}
\newcommand{\srb}[1]{\textcolor{purple}{SrB: #1}}

\ExplSyntaxOn
\cs_gset:Npn \__first_footerline:
  { \group_begin: \small \sffamily \__short_authors: \group_end: }
\ExplSyntaxOff

\begin{document}

\let\WriteBookmarks\relax
\def\floatpagepagefraction{1}
\def\textpagefraction{.001}

\shorttitle{Seymour second neighbourhood conjecture when $\delta=7$}
\shortauthors{Sadhukhan et al.}

\title[mode=title]{A proof of Seymour's second neighborhood conjecture for oriented graphs with minimum out-degree equal to \(7\)}


\author[1]{Arpan Sadhukhan}
\ead{ra.a.sadhukhan@iitdh.ac.in}

\author[2]{R. B. Sandeep}
%\cormark[1]
\ead{sandeeprb@iitdh.ac.in}

\author[1]{Sagnik Sen}
\ead{sen@iitdh.ac.in}

\affiliation[1]{
  organization={Department of Mathematics, Indian Institute of Technology Dharwad},
  %addressline={Address line},
  %city={Dharwad},
  %postcode={Postcode},
  %state={State},
  %country={India}
}

\affiliation[2]{
  organization={Department of Computer Science and Engineering, Indian Institute of Technology Dharwad},
  %addressline={Address line},
  %city={Dharwad},
  %postcode={Postcode},
  %state={State},
  %country={India}
}



%\cortext[1]{Corresponding author}

\begin{abstract}
We prove Seymour's second neighborhood conjecture on oriented graphs whose minimum out-degree is equal to $7$. This gives, to our knowledge, the first improvement of the minimum out-degree threshold in two decades, since the work of Kaneko and Locke in 2001, who resolved the conjecture for oriented graphs whose minimum out-degree is at most $6$. The proof is partially computer-assisted: after a sequence of local reductions, the remaining finite obstruction models are eliminated by reproducible OR-Tools CP-SAT infeasibility checks. %CP-SAT infeasibility checks and audited certificate logs.
\end{abstract}

\begin{keywords}
Seymour second neighbourhood \sep minimum out-degree  
\end{keywords}

\maketitle

\section{Introduction}

Let \(D\) be an oriented graph which denotes a finite directed simple graph without loops or digons; that is, there are no vertices $u,v$ for which the arcs $(u,v)$ and $(v,u)$ are both included. For a vertex \(v\in V(D)\), let \[ N^+(v)=\{u\in V(D): vu\in E(D)\} \] denote the first out-neighbourhood of \(v\), and \[ N^{++}(v)=\{w\in V(D)\setminus N^+(v): \text{ there exists }u\in N^+(v)\text{ with }uw\in E(D)\} \] for the second out-neighbourhood of \(v\). A vertex \(v\) is called a \emph{Seymour vertex} if \[ |N^{++}(v)|\ge |N^+(v)|. \]

In the early 1990s, Seymour posed the now-celebrated second neighborhood conjecture: every oriented graph has a Seymour vertex. Seymour's second neighbourhood conjecture was first published by Dean and Latka in connection with tournaments~\cite{DeanLatka1995}. The tournament case, known as Dean's conjecture, was proved by Fisher~\cite{Fisher1996}; Havet and Thomassé later gave a different proof using median orders~\cite{HavetThomasse2000}. The conjecture has since attracted substantial attention; see, for example,~\cite{AiGerkeGutinWangYeoZhou2024,BrantnerBrockmanKaySnively2009, CohnGodboleHarknessZhang2016,DaraFrancisJacobNarayanan2022, EspunyDiazGiraoGranetKronenberg2024,FidlerYuster2007,Ghazal2012, HernandezCruzGaleanaSanchez2012,KanekoLocke2001,LiangXu2017, Lim2020,Llado2013,MezherDaamouch2024,Seacrest2018, guo2026seymourtightorientations, huang2024improvedboundseymourssecond}. The conjecture has remained largely open and solved for very specific cases. In particular, Kaneko and Locke~\cite{KanekoLocke2001} proved the conjecture for every oriented graph with a vertex of out-degree at most \(6\) in 2001. This threshold remained unchanged since then.

We improve the known minimum outdegree bound from \(6\) to \(7\), proving the conjecture for every oriented graph with minimum out-degree at most \(7\). In particular, we prove the following stronger version:
\begin{theorem}
    Let $D$ be an oriented graph with minimum out-degree less than or equal to $7$. Let $s\in V(D)$ be any vertex of minimum out-degree. Then, there exists a seymour vertex in $s\cup N^+(s) \cup N^{++}(s)$.
\end{theorem}
To the best of our knowledge, this result is the first improvement in more than two decades on the minimum-out-degree approach to Seymour's Second Neighborhood Conjecture. The organization of the paper, together with a brief overview of the proof, is given in the preliminaries below.

\section{Preliminaries}
Throughout, \(D\) is a finite simple oriented graph with first and second out-neighborhood as defined in the introduction. Let $d^+(v)=|N^{+}(v)|$ denote the out-degree of $v$. We say that a vertex $v$ \emph{out-dominates} a vertex $u$ if $u\in N^+(v)$.
%Let $\delta^+(D)$ denote minimum out-degree of a vertex in $D$. 
For a set \(X\subseteq V(D)\), define
\[
N^{+}(X)=\bigcup_{x\in X}N^{+}(x).
\]

If $a\in X, \; d_X^+(a)$ denote the number of out-neighbors of $a$ in $X$. The out-distance from \(v\) to \(w\) denoted by $d^+(v,w)$ is the length of a shortest directed path from \(v\) to \(w\), if such a path exists. For vertex sets $X,Y$, define
\[
e(X,Y)=|\{(x,y)\in X\times Y:x\to y\}|, \; \; \text{where $x\to y$ denotes that $y$ is an out-neighbor of $x$}.
\]




If there is an arc (directed edge) from $x$ to $y$, then $x$ is denoted as the head and $y$ is denoted as the tail of the arc. With slight abuse of notation, edges refers to arcs throughout the paper. In this paper, we denote by \(s\) a fixed vertex of minimum out-degree of the oriented graph $D$ that we work with, with slight abuse of notation we denote $d^+(s)$ by both $\delta^+(D)$ and $\delta$.


\begin{definition}
   We denote, \(A=N^{+}(s),\;
B=N^{+}(A)\setminus A \; \text{and} \;
C=N^{+}(B)\setminus(A\cup B).
\) 
\end{definition}
Then $|A|=\delta$, every vertex of \(A\) is at out-distance \(1\) from \(s\), and every vertex of \(B\) is at out-distance \(2\) from \(s\).

\begin{definition}
   Let \(a_{1}\in A\), denote a vertex of $A$ such that \(|N^{+}(a_{1})\cap A|\) is minimum. If several vertices of \(A\) attain this minimum, break the tie by choosing one for which \(|N^{+}(a_{1})\cap B|\) is minimum. 
\end{definition} 

\begin{definition}Define, \(
A_{1}=N^{+}(a_{1})\cap A, \; \text{and} \;
A_{2}=N^{+}(A_{1})\cap (A\setminus A_{1}).
\)
    
\end{definition}

For a vertex \(a\in A\), all out-neighbours of \(a\) lie in \(A\cup B\), thus, $d^+(a_{1})=|N^{+}(a_{1})\cap A|+|N^{+}(a_{1})\cap B|
=|A_{1}|+|N^{+}(a_{1})\cap B|.$ \\

Now we will state two lemmas and a proposition that are already present in the paper~\cite{KanekoLocke2001}. We give the proofs in the appendix~\ref{sec:Missing proofs in preliminaries} for completeness. These two lemmas and the proposition will help us divide the problem into relevant cases. First, we state the general upper bound on the size of  $A_1$ with respect to the minimum out-degree of $D$.

\begin{lemma}\label{lem:general bound for A_1}
\[
|A_{1}|\leq \left\lceil \frac{\delta}{2}\right\rceil-1.
\]
\end{lemma}

\begin{proof}
See Appendix~\ref{app:lem:general bound for A_1}
\end{proof}

We shall also use the following auxiliary fact,  from the paper~\cite{KanekoLocke2001} which will be used often in the paper.

\begin{lemma}~\label{lem:perfectdegree}
Suppose that \(|A|\) is odd and
\[
|A_{1}|=\frac{|A|-1}{2}.
\]
Then D[A] is a regular tournament and $|A_2|=|A_{1}|.$
\end{lemma}

\begin{proof}
See Appendix~\ref{app:lem:perfectdegree}.

\end{proof}

Next we note the following proposition from~\cite{KanekoLocke2001} that shows existence of Seymour vertex under certain relaxed conditions. 

\begin{proposition}~\label{prop:easyseymour}

\begin{enumerate}
    \item If $|B|\geq \delta$, then $s$ is a Seymour vertex.

\item If $|B|\leq \left\lceil \frac{\delta}{2}\right\rceil$
and \(\delta\geq 2\), then $a_1$ is a Seymour vertex.
\end{enumerate}
\end{proposition}

\begin{proof}
See Appendix~\ref{app: prop:easyseymour}.
\end{proof}

\subsection*{Outline of the paper}
In the rest of the paper we consider oriented graphs with minimum out-degree exactly equal to 7. Recall that $s$ denotes a vertex with minimum out-degree. Thus if $B\geq 7$ then clearly $s$ is a Seymour vertex. Also, by the above proposition if $B \leq 4$, then $a_1$ is a Seymour vertex. Thus in order to resolve Seymour's conjecture for oriented graphs with minimum out-degree equal to 7, we only need to show the existence of a Seymour vertex in the following cases.

\begin{enumerate}
    \item $|A|=7, |B|=5$
    \item $|A|=7, |B|=6$
\end{enumerate}

We divide the second case above into 3 sub-cases based on the size of $A_1$. By lemma~\ref{lem:general bound for A_1}, we know $|A_1| \leq 3$. Also, $A_1\geq 1$, otherwise the minimum out-degree requirement of $a_1$ is not satisfied when $|B|=6$. Thus the only cases where we need to show the existence of a Seymour vertex are the following.

\begin{enumerate}
    \item $|A|=7, |B|=5$
    \item $|A|=7, |B|=6, |A_1|=3$
    \item $|A|=7, |B|=6, |A_1|=2$
    \item $|A|=7, |B|=6, |A_1|=1$
\end{enumerate}

Each of the 4 cases above is handled in separate sections below. The proof of existence of a Seymour vertex in cases 1 and 2 are manual and requires no computer assistance, whereas the parts of the proof of existence of a Seymour vertex in cases 3 and 4 is computer assisted after effectively reducing the possibilities to a finite number of configurations based on certain assumptions. The rest of the paper following the overview of the proof is divided into four sections, with each section containing the proof of the existence of a Seymour vertex in each of the above mentioned cases and thus proving our claim.

\subsection*{Overview of the proof}
The first two cases, when $|B|=5$ and when $|B|=6, |A_1|=3$ are relatively easier than the rest. It only relies on good structural observation in a very local neighborhood of the smallest degree vertex, mostly using counting arguments. The case $|B|=6, |A_1|=2$ and the case $|B|=6, |A_1|=1$ are surprisingly much harder to resolve. The reason behind it is that the conditions in the first two cases like $|A_1|=3$, converts the whole problem into a near-equality situation where almost every counting inequality is tight. Whereas like when $|A_1|=2$, it leaves too much slack, so the same counts stop forcing structure.
There we first analyze the local structure determined by a minimum out-degree vertex \(s\), with the aim of proving that \(A\cup B\) already contains a Seymour vertex. The first step is to study the controlled part \(A\cup B\cup C\). From this local structure we obtain deterministic constraints, together with finite ranges for the parameters that can occur.

The key idea is that we do not need to know the entire graph outside \(A\cup B\cup C\) in order to decide whether a vertex in \(A\cup B\) is Seymour. We only need to know which outside vertices can appear as second-neighbors of vertices in \(A\cup B\). Thus the part of the graph beyond \(A\cup B\cup C\) can be compressed. More precisely, such an outside vertex matters only through the set of vertices of \(C\) that send an edge to it. We call this set its signature. Vertices with the same signature contribute in exactly the same way to the second-neighborhood counts of vertices in \(A\cup B\), so they can be grouped together. In this way, the relevant information about the outside world is captured by finitely many signature variables and the non-Seymour inequalities for the vertices in \(A\cup B\).

Something we call the 'Trimming lemma' is then used to justify this reduction. It is a very simple lemma that essentially says that if we delete edges whose tails lie outside \(A\cup B\), then the first out-neighborhood of every vertex in \(A\cup B\) is unchanged, while its second out-neighborhood can only become smaller. Therefore, if all vertices in \(A\cup B\) were non-Seymour before the deletion, they remain non-Seymour after the deletion. This allows us to delete irrelevant outgoing edges from vertices outside \(A\cup B\). After this deletion, the remaining outside behavior is completely described by finitely many signature variables, whose ranges are made bounded.

Consequently, if there were a counterexample in which every vertex of \(A\cup B\) is non-Seymour, then the deterministic constraints and the non-Seymour inequalities would have a feasible solution in one of the finite parameter ranges obtained manually. The computer-assisted part checks precisely these finite systems and finds that all of them are infeasible. Hence no such local obstruction exists, and therefore \(A\cup B\) contains a Seymour vertex.

The CP-SAT models used in the computational part were developed with assistance from ChatGPT 5.5 Pro; the resulting OR-Tools CP-SAT encodings were then run and independently audited by the authors to verify that all corresponding finite configurations are infeasible. The files can be found at \url{https://github.com/rbsandeep/Seymour-Vertex-delta7}.







\section{The case  $|A|=7$ and $|B|=5$}

In this section we resolve Seymour’s second neighborhood conjecture when $|A|=7$ and $|B|=5$. The proof is based on counting and logical arguments. No computer search or certificate is used in this case. 

Recall that $s$ is a minimum out-degree vertex with $d^+(s)=7$, with $A=N^+(s),\; B=N^{++}(s)$. Then $|A|=7$. Since $s\to a$ for every $a\in A$, $N^+(a)\subseteq A\cup B$.



Recall that, $C=N^+(B)\setminus (A\cup B).$
Note that the set $C$ may contain the root $s$ if some vertex of $B$ sends to $s$. The only property needed is that every out-neighbor of a vertex of $B$ lies in $A\cup B\cup C$.
In particular, we prove the following theorem below.

\begin{theorem}
Let $D$ be a finite oriented graph with $\delta^+(D)=7$. Let $s$ be a vertex with $d^+(s)=7$, put $A=N^+(s)$ and $B=N^{++}(s)$. If $|B|=5$, then some vertex of $A$ is a Seymour vertex.
\end{theorem}

\begin{proof}
We know that every out-neighbor of a vertex of $A$ lies in $A\cup B$. Since every vertex of $D$ has outdegree at least $7$ and $|B|=5$, each vertex of $A$ has at least two out-neighbors in $A$, that is, $|N^+(a)\cap A|\ge 7-|B|=2$ for all $a\in A$.

Recall that $a_1\in A$ is the vertex of $A$ that minimizes $|N^+(a)\cap A|$ over $a\in A$, and put and $A_1=N^+(a_1)\cap A$.
Since $e(D[A])\leq \binom{7}{2}=21$, it is easy to see $|A_1|\leq 3$. Thus $|A_1|\in\{2,3\}$.

We will now handle the two cases separately.

\medskip
\noindent\emph{Case 1: $|A_1|=2$.}


Since $N^+(a_1)\subseteq A\cup B$, $|A_1|=2$, and $|B|=5$, the set $A_1\cup B$ has exactly seven vertices. Now out-degree of $a_1$ is at least seven, thus $N^+(a_1)=A_1\cup B$ or in other words $a_1$ out-dominates all vertices of $B$ and $d^+(a_1)=7$.

Recall that $A_2=(N^+(A_1)\cap A)\setminus A_1.$ By the minimality of $a_1$, each vertex of $A_1$ has at least two out-neighbors in $A$, so the two vertices of $A_1$ send at least four edges into $A$. At most one of these edges can stay inside $A_1$. Therefore at least three edges go from $A_1$ to $A\setminus (A_1\cup\{a_1\})$. A fixed vertex can receive at most two of these edges from the two vertices of $A_1$, hence $|A_2|\geq 2$.


Since $a_1$ out-dominates every vertex of $B$, every vertex of $C$ is a second out-neighbor of $a_1$, and every vertex of $A_2$ is also a second out-neighbor of $a_1$. The sets $A_2$ and $C$ are disjoint and neither is contained in $N^+(a_1)$. Thus, if $|C|\ge 5$, then,
\[
|N^{++}(a_1)|\ge |A_2|+|C|\ge 2+5=7=d^+(a_1),
\]
so $a_1$ is Seymour. Hence for the sake of contradiction assume $ |C|\le 4.$
We now count edges between the layers. Every vertex of $A$ has out-degree at least $7$ and all its out-neighbors lie in $A\cup B$, so
\[
e(A,A)+e(A,B)\ge 7\cdot 7=49.
\]
Using $e(A,A)\le 21$, we get $e(A,B)\ge 28.$ There are $7\cdot 5=35$ cross-pairs between $A$ and $B$, thus, $e(B,A)\le 35-28=7.$
The vertices of $B$ have total out-degree at least $5\cdot 7=35$, and their out-neighbors lie in $A\cup B\cup C$. Since $e(B,B)\le \binom{5}{2}=10$, We have,
$e(B,C)\ge 35-7-10=18$.

Each vertex of $C$ receives at most five edges from $B$, hence $|C|\ge \lceil 18/5\rceil=4$. Thus we may assume $C=4$.





Thus, $e(B,C)\le 5|C|=20$ and $e(B,B)\le 10$, thus by the inequality $e(B,A)+e(B,B)+e(B,C)\ge 35$, we have, $e(B,A)\ge 35-10-20=5.$
If $|A_2|\ge 3$, then clearly, $|N^{++}(a_1)|\ge |A_2|+|C|\ge 3+4=7=d^+(a_1)$,
so $a_1$ is Seymour. Thus, we may assume $|A_2|=2.$

Let $A_3=A\setminus (\{a_1\}\cup A_1\cup A_2).$ Then $|A_3|=2$. 

If some vertex of $B$ sent an edge to a vertex of $A_3$, that vertex of $A_3$ would be a second out-neighbor of $a_1$ through $B$, and then
\[
|N^{++}(a_1)|\ge |A_2|+|C|+1=2+4+1=7=d^+(a_1).
\]
Thus, if $a_1$ is not Seymour, no edge goes from $B$ to $A_3$. Also, no edge goes from $B$ to $a_1$, because $a_1$ out-dominates all vertices of $B$. Hence, every edge from $B$ to $A$ lands in $U=A_1\cup A_2.$ We also have $e(B,A)= e(B,U)\ge 5$. Now we will derive a contradiction by showing this cannot happen.


Let $m=e(A_1,A_2)$. Clearly from definition it follows that $m\geq 2$. The edges from $A_1$ to $A$ can only go inside $A_1$ or to $A_2$: they cannot go to $a_1$, and they cannot go to $A_3$ by the definition of $A_2$. Since at most one edge lies inside $A_1$, $e(A_1,A)\le m+1$.
Next count edges from $A_2$ into $A$. There are at most $2$ edges from $A_2$ to $a_1$, at most $4-m$ edges from $A_2$ to $A_1$, at most one edge inside $A_2$, and at most $|A_2||A_3|=4$ edges from $A_2$ to $A_3$. Therefore
\[
 e(A_2,A)\le 2+(4-m)+1+4=11-m.
\]
Adding together we have, $e(U,A)\le 12$.
The four vertices of $U$ have total outdegree at least $4\cdot 7=28$, and all their out-neighbors lie in $A\cup B$. Hence
\[
 e(U,B)\ge 28-e(U,A)\ge 16.
\]
Between $U$ and $B$ there are only $4\cdot 5=20$ cross-pairs. By orientation and (23),
\[
 e(B,U)\le 20-16=4,
\]
which is a contradiction. This finishes the proof for \emph{Case 1}.

\medskip
\noindent\emph{Case 2: $|A_1|=3$.}
Since $a_1$ minimizes $|N^+(a)\cap A|$ and $|A_1|=3$, by Lemma~\ref{lem:perfectdegree}, we know that $D[A]$ is a tournament and for every vertex $a\in A$, $|N^+(a) \cap A|=3$ and also $|N^{++}(a)\cap A|=3$. Now since all vertex of $A$ has the same out-degree inside $A$, by definition of $a_1$, we choose it in such a way that $|N^+(a)\cap B|$ is minimized when $a=a_1$.

Now observe that if $|N^+(a_1)\cap B|=5$ then $|N^+(a)\cap B|=5$ for all $a\in A$ which implies $e(B,A)=0$, thus $e(B,C)\geq 35-10=25$, thus $|C|\geq 5$. Also, $C\subseteq N^{++}(a_1)$ and $C\cap A=\emptyset$. Thus $|N^{++}(a_1)|=|N^{++}(a_1)\cap C| + |N^{++}(a_1)\cap A|\geq 5+3=8=|N^+(a_1)|$. Thus, $a_1$ is a Seymour vertex. Hence, we may assume $|N^{+}(a_1)\cap B|=4$ (minimum degree of $a_1$ being greater than or equal to $7$ forces the equality here).

Also, from the analysis of \emph{Case 1}, we know that $e(B,C)\geq 18$ (this was derived independent of the size of $A_1$ thus it holds in this case too). Let $E=N^{+}(a_1)\cap B$. Now, observe that, $|N^+(E)\cap C|\leq 3$, otherwise $a_1$ would be Seymour vertex. Thus, $e(E,C)\leq 12$ (as $|E|=4$). Now $e(B,C)=e(E,C)+e(B\setminus E,C)$. Thus, $e(B\setminus E,C)\geq 6$. Now, $B\setminus E$ consists of a singleton vertex, say $u$. Clearly, there exists $a^*\in A$, such that $u\in N^+(a^*)$. Thus $|N^{++}(a^*)|=|N^{++}(a^*)\cap C| + |N^{++}(a^*)\cap A|\geq 3+6=9$. Also, $|N^+(a^*)|\leq 8$. Thus, $a^*$ is a Seymour vertex. This finishes the proof of \emph{Case 2}.




\end{proof}


\section{Case $|A|=7, |B|=6, |A_1|=3$}
In this section we resolve Seymour’s second neighborhood conjecture when $|A| = 7, |B| = 6, |A_1|=3$. The proof is based on counting and logical arguments. No computer search or certificate is used in this case.


\begin{theorem}~\label{thm:Case B6A_1_3}
Let $D$ be a finite oriented graph with $\delta^+(D)=7$. Let $s$ be a vertex with $d^+(s)=7$, put $A=N^+(s)$ and $B=N^{++}(s)$, and assume $|B|=6$. Choose $a_1\in A$ minimizing $d_A^+(a)$ over $a\in A$, and set $A_1=N_A^+(a_1)$. If $|A_1|=3$, then $D$ has a Seymour vertex in $A\cup B$.
\end{theorem}


As a direct consequence of Lemma~\ref{lem:perfectdegree} we know that $D[A]$ is a tournament and for every vertex $a\in A$, $|N^+(a) \cap A|=3$ and also $|N^{++}(a)\cap A|=3$. Now since all vertex of $A$ has same out-degree inside $A$, by definition of $a_1$, we choose it in such a way that $|N^+(a)\cap B|$ is minimized when $a=a_1$.
The proof splits according to how many vertices of $B$ are out-dominated by $a_1$. Let $P=N^+(a_1)\cap B$ and $Q=B\setminus P$. It is clear that  $|P| \in \{4,5,6\}$ as $|B|=6$ and $|N^+(a_1)|\geq 7$.



%For $a\in A$, define
%\[
%p(a)=d_B^+(a)=|N^+(a)\cap B|.
%\]
%Clearly, $d^+(a)=3+p(a)$.
%Since every vertex has out-degree at least $7$ and $|B|=6$,
%\begin{equation}
%p(a)\in \{4,5,6\}\qquad (a\in A). \tag{3}
%\end{equation}
%It remains to analyze these three possibilities.

%For the rest of the proof define
%\[
%C=N^+(B)\setminus (A\cup B).
%\]
%This set may contain the root $s$ if some vertex of $B$ sends to $s$. This causes no difficulty. The only property used below is that every out-neighbor of a vertex of $B$ lies in $A\cup B\cup C$.

\subsection*{A vertex of $A$ has four out-neighbors in $B$}
This is the case in which $|P|=4$, $|Q|=2$ and $d^+(a_1)=3+4=7$ as $|N^+(a_1)|=|N^+(a_1)\cap A|+|P|$.

Suppose, for contradiction, that $a_1$ is not Seymour. We know that $a_1$ has exactly three second out-neighbors in $A$. Since $d^+(a_1)=7$, we may assume the number of second out-neighbors of $a$ outside $A$ is at most $3$. Let $Q=\{b_5,b_6\}$. Let $C_P=N^+(P)\cap C$. We now further divide this case into sub-cases depending on how many vertices of $Q$ are there in the second out-neighborhood of $a_1$.
\medskip
%Every vertex of $T\cup C_P$ is a second out-neighbor of $a$ outside $A$, and the two sets are disjoint. Hence
%\begin{equation}
%q+|C_P|\le 3. \tag{4}
%\end{equation}

\emph{Sub-case: $|Q\cap N^{++}(a_1)|=0$}


 No vertex of $Q$ is a second out-neighbor of $a_1$. Therefore no vertex of $A_1 \cup P$ sends an edge to $Q$. Also, $a$ itself sends no edge to $Q$. Hence the only possible edges from $A$ to $Q$ have their tails in $A_2$ which has size equal to $3$. Thus, $e(A,Q)\le 3\cdot 2=6$.
Now as every vertex of $A$ sends at least four edges to $B$, we have $e(A,B)\ge 7\cdot 4=28$. Thus, $e(A,P)=e(A,B)-e(A,Q)\ge 28-6=22$.


Between $A$ and $P$ there are $7\cdot 4=28$ arcs. Therefore, $e(P,A)\le 28-22=6$. 

The four vertices of $P$ have total outdegree at least $28$. They send at most $6$ edges to $A$ by (6), at most $\binom{4}{2}=6$ edges inside $P$, and no edges to $Q$. Hence $e(P,C)\ge 28-6-6=16$.
All these edges land in $C_P$, and a vertex of $C_P$ receives at most four edges from $P$. Hence $|C_P|\ge \lceil 16/4\rceil=4$. Thus, $|N^{++}(a_1)|=|A_2|+|C_P|\geq 3+4=7=|N^+(a_1)|$, thus $a_1$ is a Seymour vertex.

\medskip
\emph{Sub-case: $|Q\cap N^{++}(a_1)|=1$}

Only one vertex of $Q$ is a second out-neighbor of $a$. Thus, $|N^{++}(a_1)|=|A_2|+|C_P|+1$. 
Now observe that, the three vertices of $A_1$ can send edges only to that one vertex of $Q$, while the three vertices of $A_2$ can send edges to both vertices of $Q$. Again, $a$ sends no edge to $Q$. Therefore, $e(A,Q)\le 3\cdot 1+3\cdot 2=9$.
As before, $e(A,P)\ge 28-9=19$, thus $e(P,A)\le 28-19=9$. The vertices of $P$ send at most $\binom{4}{2}=6$ edges inside $P$, and at most $4$ edges to $Q$, because all such edges must land in the unique vertex of $Q\cap N^{++}(a)$. Counting the out-degrees of the four vertices of $P$ gives $e(P,C)\ge 28-9-6-4=9$.
Therefore $|C_P|\ge \lceil 9/4\rceil=3$. Thus, $a_1$ is a Seymour vertex.

\medskip
 \emph{Sub-case: $|Q\cap N^{++}(a_1)|=2$}

Here both vertices of $Q$ are second out-neighbors of $a_1$. Thus, $|N^{++}(a_1)|=|A_2|+|C_P|+2$. Hence if $C_p\geq 2$, $a_1$ is a Seymour vertex.

Now observe that, the three vertices of $A_1$ and that of $A_2$ all can send edges to both vertices of $Q$. $e(A,Q)\le 3\cdot 2+3\cdot 2=12$.

As before, $e(A,P)\ge 28-12=16$, thus $e(P,A)\le 28-16=12$. The vertices of $P$ send at most $\binom{4}{2}=6$ edges inside $P$, and at most $8$ edges to $Q$, Counting the out-degrees of the four vertices of $P$ gives $e(P,C)\ge 28-12-6-8=2$. Therefore $|C_P|\ge \lceil 2/4\rceil=1$. Therefore we may assume $|C_P|=1$, otherwise we are done.

Recall that as every vertex of $A$ sends at least four edges to $B$, we have $e(A,B)\ge 7\cdot 4=28$, thus $e(B,A)\leq 42-28=14$. Now $e(B,C)\geq 7|B|-\binom{6}{2}-e(B,A)$, thus $e(B,C)\geq 13$. Now since $|C_P|=1$, we have, $e(P,C)\leq 4$. Thus, $e(Q,C)\geq 9$. Hence, either $e(b_5,C)\geq 5$ or $e(b_6,C)\geq 5$. Without loss of generality assume, $e(b_6,C)\geq 5$. Let $a^*\in A$ be the vertex such that $b_6\in N^+(a^*)$. Thus $|N^{++}(a^*)|\geq |N^+(b_6)\cap C|+|N^{++}(a^*)\cap A|\geq 5+3=8$. Thus, $a^*$ is not a Seymour vertex iff $N^+(a^*)\cap B=B$.

Now if $N^+(a^*)\cap B=B$, we can refine the earlier calculations and now we have $e(A,B)\geq 28+2=30$, thus $e(B,A)\leq 42-30=12$, thus $e(B,C)\geq 42-15-12=15$. Thus $e(Q,C)\geq 11$. Thus, there exists a vertex $x\in \{b_5, b_6\}$ in Q with $e(x,C)\geq 6$, hence the vertex of $A$ that out-dominates $x$ must be a Seymour vertex. This finishes the case.



\subsection*{Every vertex of $A$ out-dominates at least $5$ vertices in $B$}
Now, we have, $N^+(a_1)\cap B=5$. Recall, $P= N^+(a_1)\cap B$ and $Q=B\setminus P$ and $C_P=N^+(P)\cap C$. Here $Q$ contains only a singleton vertex, call it $q$. Also, since $N^+(a_1)\cap B=5$, we have, $N^+(a)\cap B\geq 5$ for all $a\in A$. Thus, $e(A,B)\geq 7\times 5=35$.



\medskip
\emph{Sub-case: $q\notin N^{++}(a_1)$.}

Observe that no vertex of $(N^+(a_1)\cap A)\cup P$ sends an edge to $q$, and $a_1$ itself does not send to $q$. Thus, the only possible edges of $A$ to $q$ have their tails in the three vertices of $A_2$. Thus, $e(A,\{q\})\le 3$. 
Therefore, $e(A,P)\ge e(A,B)-3=32$. Hence, $e(P,A)\le 35-32=3$. The five vertices of $P$ have total out-degree at least $35$. They send at most $3$ edges to $A$, at most $\binom{5}{2}=10$ edges inside $P$, and no edges to $q$. Therefore, $e(P,C)\ge 35-3-10=22$. Thus, $C_P \geq \lceil 22/5\rceil=5$. Hence, $|N^{++}(a_1)|=|N^{++}(a_1)\cap A|$+$|N^{++}(a_1)\cap C|=|A_2|+|C_P|\geq 3+5=8=|N^+(a_1)|$. Hence $a_1$ is a Seymour vertex.

\medskip
\emph{Sub-case: $q\in N^{++}(a_1)$.}

Observe that, $|N^{++}(a_1)|=|A_2|+|C_P|+1$, the $+1$ is present because of the vertex $q$, also $|N^+(a_1)|=8$. Now if $e(P,C_P)\geq 16$ then $|C_P|\geq 4$, thus $a_1$ is a Seymour vertex. Thus, we may assume, $|C_P|\leq 3$ and $e(P,C_P)\leq 15$.

Now suppose $e(q,C)\geq 6$. Let $a^*\in A$ be such that $q\in N^+(a^*)$. Then the first out-neighborhood of $a^*$ has size at most $9$ and the second out-neighborhood has size at least $9$, hence $a^*$ will be a Seymour vertex. Thus, we may assume $e(q,C)\leq 5$. 

Combining the two above we may assume $e(B,C)\leq 20$.

Now let $t=e(A,\{q\})$ and $u=e(P,\{q\})$. Again $e(A,B)\ge 35$, and therefore
$e(A,P)\ge 35-t$, and $e(P,A)\le t.$
Counting the outdegrees of the five vertices of $P$ gives $e(P,C)\ge 35-e(P,A)-e(P,P)-e(P,\{q\})\ge 35-t-10-u=25-t-u$.

We now bound the out-degree of $q$ inside $A\cup B$. The vertex $q$ may send an edge to $a$, because $a$ does not send to $q$. Among the other six vertices of $A$, exactly $t$ send edges to $q$, so $q$ can send to at most $6-t$ of them. Among the five vertices of $P$, exactly $u$ send edges to $q$, so $q$ can send to at most $5-u$ of them. Thus,
\[
|N^+(q)\cap (A\cup B)|\le 1+(6-t)+(5-u)=12-t-u.
\]
Thus, $e(q,C)\geq 7-(12-t-u)=t+u-5$. Combining, we have,
 \[
 e(B,C)\geq e(P,C)+e(q,C)\geq 25-t-u+t+u-5=20
 .\]
From the previous assumption, this forces an equality, thus $e(B,C)=20, e(P,C)=15$ and $e(q,C)=5$, otherwise we have a Seymour vertex. 
Thus any vertex of $A$ whose out-neighbor is $q$ must out-dominate whole of $B$, otherwise it will be a Seymour vertex. Thus, $e(A,B)\geq 35+t$. Also, $t\geq 1$ from definition. Hence, $e(B,A)\leq 42-(35+t)=7-t$. Hence,
\[
e(B,C)\ge 42-(7-t)-15=20+t\geq 21
\]
which is a contradiction. This finishes the proof.



\subsection*{Every vertex of $A$ out-dominates $B$}

Now we are in the case where every vertex of $A$ out-dominates all vertices in $B$. Thus $e(B,A)=0$, hence $e(B,C)\geq 42-\binom{6}{2}=27$. Thus, $|C|\geq 5$. Also, if $|C|\geq 6$, then all the vertices of $A$ is a Seymour vertex as $A_2\cup C \subseteq N^{++}(a)$ and $N^+(a)\leq 9$ for all $a\in A$. Thus, we may assume $|C|=5$.



For every $b\in B$, all out-neighbors of $b$ lie in $B\cup C$, so $|N^+(b)\cap B|+|N^+(b)\cap B|=d^+(b)\ge 7$. Since $|C|=5$, this implies $|N^+(b)\cap B|\ge 2$ for all \(b\in B\). 


Choose $b\in B$ with $d_B^+(b)$ minimum. Since $e(B,B)\le 15$, the average out-degree of a vertex of $B$ is at most $15/6<3$, this implies $|N^+(b)\cap B|=2$.

Since, $|C|=5$, we also have $|N^+(b)\cap C|=5$, and $d^+(b)=7$. Thus $b$ dominates every vertex of $C$.

Define $B_1=N^+(b)\cap B$, thus $|B_1|=2$ and $Y=B\setminus (\{b\}\cup B_1)$, thus $|Y|=3$. Also, define $S=N^+(B_1)\cap Y$.


\begin{claim}
$|S|\ge 2$.
\end{claim}

\begin{proof}
By definition, $e(B_1, b)=0$. Now each vertex of $B_1$ has at least two out-neighbors in $B$. Thus, $e(B_1,B)\geq 4$. Now as $|B_1|=2$, $e(B_1,B_1)\leq 1$. So $e(B_1, B\setminus (B_1\cup \{b\}))\geq 3$. Now if $|S|\leq 1$, then clearly, $e(B_1, B\setminus (B_1\cup \{b\})\leq 2$, which is a contradiction.
\end{proof}



Define $R=N^+(C)\setminus (C\cup B)$.
Observe that $e(B,C)\geq 27$ and $e(b,C)=5$, thus we have $e(B\setminus \{b\},C)\geq 22$, thus $e(C,B)=e(C,B\setminus \{b\})\leq 5\times 5-22=3$. Now, since each vertex of $C$ has out-degree at least $7$ and there are at most $10$ edges inside $C$, we have $e(C,R)\geq 35-10-3=22$. Thus, $|R|\ge \left\lceil\frac{22}{5}\right\rceil=5$. Also, clearly, $R\in N^{++}(b)\setminus B$. Thus $|N^{++}(b)|\geq |R|+ |S|\geq 5+2=7=N^+(b)$. Thus $b$ is a Seymour vertex. This finishes the proof of this case.



\begin{proof}[Proof of Theorem~\ref{thm:Case B6A_1_3}]

By Lemma 1, every vertex $a\in A$ has $d_A^+(a)=3$, and therefore $|P|\in \{4,5,6\}$. We have successfully shown existence of a Seymour vertex in all the three cases. This finishes the proof.
\end{proof}

\section{Case  $|A|=7, |B|=6, |A_1|=2$}\label{sec:B6A12}
In this section we resolve Seymour’s second neighborhood conjecture when $|A| = 7, |B| = 6, |A_1|=2$. In particular, we prove the following theorem:

\begin{theorem}\label{thm:main}
Let $D$ be a finite oriented graph with $\delta^+(D)=7$.  Let $s$ be a vertex with $d^+(s)=7$, and put $A=N^+(s)$ and $B=N^{++}(s)$.
Assume $|B|=6$ and $\min_{a\in A}|N^+(a)\cap A|=2$.
Then some vertex in $A\cup B$ is a Seymour vertex.
\end{theorem}


Write $r=|N^+(a_1)\cap B|.$
Now, every out-neighbor of $a_1$ lies in $A\cup B$. Since $\delta^+(D)=7$, we have $d^+(a_1)\ge 7$, so $2+r\ge 7$. Since $|B|=6$, the only possibilities are $r=5$ and $r=6.$
The proof treats these two branches separately. We first give a general bound. Let $A_1=\{u_1,u_2\}$.

\begin{lemma}\label{lem:X-at-least-2}
Let $A_2=X_0=(N^+(A_1)\cap A)\setminus A_1$. Then $|X_0|\ge 2$.
\end{lemma}
\begin{proof}
Each of $u_1,u_2$ has at least two out-neighbors in $A$, because $a_1$ was chosen with the minimum possible value $|N^+(a)\cap A|=2$.  Neither $u_i$ sends to $a_1$, because $a_1\to u_i$ and the graph is oriented.  Among $u_1,u_2$, at most one directed edge is present, so at most one of the required edges from $\{u_1,u_2\}$ into $A$ can end in $A_1$.

Thus the at least four edges from $\{u_1,u_2\}$ into $A$ include at least three edges ending in $A\setminus(\{a_1\}\cup A_1)$.  A fixed vertex can receive at most two such edges, one from $u_1$ and one from $u_2$.  Hence at least two distinct vertices must be present in $X_0$. This finishes the proof.
\end{proof}


On a high level, the idea is to convert the problem into a problem of existence of a feasible solution of a system of finitely many boolean or integer linear constraints and inequalities with finite number of variables and whose possible values are bounded. The infeasibility of which results in a Seymour vertex in \(A\cup B\). In view of this, the following terminology that we develop in the subsection below, will be used throughout the remainder of the paper to describe the finite model. 




\subsection{Witnesses, obstructions, finite models, and trimming}
\begin{definition}[Witness set]
A \emph{witness set} is a specified set $W$ of vertices for which the non-Seymour inequalities are imposed. In this paper $W=A\cup B$.
\end{definition}

\begin{definition}[Obstruction]
For a fixed witness set $W$, an \emph{obstruction} is a configuration satisfying all standing graph-theoretic assumptions in which every vertex of $W$ is non-Seymour.
\end{definition}

\begin{definition}[Finite model]
A \emph{finite model} is a finite encoding of all information that can affect the first- and second-neighborhood inequalities for the chosen witness vertices. In this proof the finite models consist of:
\begin{enumerate}[label=(\roman*)]
\item an explicitly modeled finite local vertex set $L$;
\item Boolean variables for directed edges inside $L$;
\item Boolean variables for local second-neighbor incidences;
\item integer variables for outside vertices with identical exact predecessor signatures;
\item deterministic constraints forced by the definitions of the relevant sets;
\item the inequalities asserting that all chosen witnesses are non-Seymour.
\end{enumerate}
A model is \emph{sound} if every genuine obstruction gives a feasible assignment of the variables.
\end{definition}

Only the implication $\text{genuine obstruction}\Longrightarrow\text{feasible finite model}$
is used. The converse is not needed; the finite models may be relaxations. Therefore, infeasibility of a sound finite model rules out the corresponding genuine obstruction.

\begin{lemma}[Trimming lemma]\label{lem:trimming}
Let $W$ be a witness set in an oriented graph $D$.  Let $D'$ be a subdigraph of $D$ obtained by deleting directed edges of $D$ whose tails all lie outside $W$.  If $v\in W$ is not a Seymour vertex in $D$, then $v$ is not a Seymour vertex in $D'$.
\end{lemma}


\begin{proof}
No edge leaving $v$ is deleted, because $v\in W$ and all deleted edges have tails outside $W$.  Hence $N^+_{D'}(v)=N^+_D(v)$.  If $z\in N^{++}_{D'}(v)$, then $v\to u\to z$ in $D'$ for some $u$, and $z\notin N^+_{D'}(v)\cup\{v\}$.  The same two directed edges are present in $D$, and the first-neighborhood of $v$ is unchanged.  Therefore $z\notin N^+_D(v)\cup\{v\}$, so $z\in N^{++}_D(v)$. Thus, for every $v\in W$, we have, $N^+(v)\cap D'=N^+(v)\cap D$ and $N^{++}(v)\cap D'\subseteq N^{++}(v)\cap D$. This finishes the proof.
\end{proof}


\begin{remark}
    The lemma will be used in two ways. First, for auxiliary non-witness vertices whose outgoing edges are represented by signature variables, we trim their outgoing edges until their total out-degree is exactly seven. Second in the $r=5$ branch, we delete all outgoing edges of the exceptional vertex $q$. Since $q$ (the unique vertex in $B\setminus N^+(a_1)$) is not a witness, the trimming lemma ensures that the infeasibility of the model will still prove our result. In both uses, all first-neighborhoods of witnesses are unchanged and second-neighborhoods of witnesses can only shrink. 
\end{remark}







\subsection{The branch $r=6$}

Here we have $N^+(a_1)\cap B=B$. Thus, $N^+(a_1)=A_1\cup B$ and $|N^+(a_1)|=8$.

\begin{definition}
    Define $X=(N^+(A_1\cup B)\cap A)\setminus A_1$ and $x=|X|$. Also, define $H=A_1\cup X$ and $R=A\setminus(\{a_1\}\cup H)$.
\end{definition}

By Lemma \ref{lem:X-at-least-2}, we have, \[2\le x\le 4,
\qquad |H|=x+2,
\qquad |R|=4-x.
\]
By the definition of $X$, no vertex of $A_1\cup B$ sends an edge to $R$.  Also no vertex of $B$ sends an edge to $a_1$, because $a_1\to B$. Recall, $C=N^+(B)\setminus(A\cup B)$.
Equivalently, this is the set of vertices outside $A\cup B$ reached from $B$; in particular, $s\in C$ exactly when some vertex of $B$ has an out-edge to $s$. Let $M=|C|$.

When $s\in C$, it is treated as an explicit auxiliary non-witness vertex. This is sound because the actual root behavior $N^+(s)=A$ is one assignment allowed by the model.  Treating the elements of $C$ uniformly is a relaxation, and infeasibility of a relaxation still rules out genuine obstructions.

\begin{lemma}\label{lem:r6-a1-exact}
In the $r=6$ branch, $N^{++}(a_1)=X\cup C$. Consequently, every $r=6$ obstruction satisfies $x+M\le 7$.

\end{lemma}
\begin{proof}
The first-neighbors of $a_1$ are exactly $A_1\cup B$.  A vertex of $A$ reached in one more step from $A_1\cup B$ lies in $A_1\cup X$: it cannot be $a_1$, because $a_1$ sends to every vertex of $A_1\cup B$, and it cannot lie in $R$ by the definition of $X$.  Vertices in $A_1$ are first-neighbors of $a_1$ and are not counted as second-neighbors, so the contribution inside $A$ is exactly $X$.

Vertices of $B$ are first-neighbors of $a_1$ and are not counted.  The remaining possible vertices that are out-neighbours of some vertex of $B$ are exactly the vertices of $C$, including the root $s$ precisely when some vertex of $B$ has an out-edge to $s$.  Hence $N^{++}(a_1)=X\cup C$.  Since $a_1$ has out-degree $8$, non-Seymourness of the witness vertex $a_1$ gives $|N^{++}(a_1)|\le 7$, and therefore $x+M\le 7$.
\end{proof}

\begin{lemma}[Counting in the $r=6$ branch]\label{lem:r6-counting}
In the $r=6$ branch,
\[
\begin{array}{c|cc}
 x & e(H,B)\ge & e(B,C)\ge\\ \hline
 2&16&19\\
 3&19&16\\
 4&23&14
\end{array}
\]
\end{lemma}
\begin{proof}
Let $h=|H|=x+2$.  Vertices of $H\subseteq A$ have all their out-neighbors in $A\cup B$ and have outdegree at least seven.  Hence the vertices of $H$ send at least $7h$ edges into $A\cup B$.

We upper-bound how many of these edges can end in $A$.  Inside $H$ there are at most $\binom{h}{2}$ directed edges.  A vertex of $X$ may send to $a_1$, giving at most $x$ further edges.  A vertex of $X$ may send to $R$, giving at most $x(4-x)$ further edges.  Vertices of $A_1$ send neither to $a_1$ nor to $R$.  Therefore
\[
e(H,B)\ge 7h-\binom{h}{2}-x-x(4-x).
\]
For $x=2,3,4$, this gives $16,19,23$.

Every vertex of $B$ has outdegree at least seven.  Vertices of $B$ cannot send to $a_1$ or to $R$.  Thus all out-neighbors of $B$ lie in $H\cup B\cup C$.  Moreover
\[
e(B,H)\le 6h-e(H,B),
\qquad e(B,B)\le \binom{6}{2}=15.
\]
Since the six vertices of $B$ have total outdegree at least $42$,
\[
e(B,C)\ge 42-(6h-e(H,B))-15.
\]
Substituting the previous lower bounds for $e(H,B)$ gives the displayed lower bounds for $e(B,C)$.
\end{proof}

Since $e(B,C)\le 6M$, Lemmas \ref{lem:r6-a1-exact} and \ref{lem:r6-counting} leave only the following rows:
\begin{equation}\label{eq:R6rows}
R_6=\{(2,4),(2,5),(3,3),(3,4),(4,3)\}.
\end{equation}
The numerical values in these rows are:
\[
\begin{array}{c|c|c|c|c|c|c}
(x,M)&|H|&|R|&|N^{++}(a_1)|=x+M&7-x-M&e(H,B)\ge&e(B,C)\ge\\ \hline
(2,4)&4&2&6&1&16&19\\
(2,5)&4&2&7&0&16&19\\
(3,3)&5&1&6&1&19&16\\
(3,4)&5&1&7&0&19&16\\
(4,3)&6&0&7&0&23&14
\end{array}
\]
In particular, in the two slack rows $(2,4)$ and $(3,3)$ the second-neighborhood of $a_1$ is still exactly $X\cup C$; the slack only means that $a_1$ has fewer than seven second-neighbors.

\subsection*{The exact $r=6$ obstruction model}\label{subsec:r6-model}
Fix a row $(x,M)\in R_6$.  Label
\[
A=\{a_1,u_1,u_2,x_1,\ldots,x_x,r_1,\ldots,r_{4-x}\},
\]
\[
A_1=\{u_1,u_2\},\qquad X=\{x_1,\ldots,x_x\},\qquad R=\{r_1,\ldots,r_{4-x}\},
\]
\[
B=\{b_1,\ldots,b_6\},
\qquad C=\{c_1,\ldots,c_M\}.
\]
The local vertex set and witness set are $L=A\cup B\cup C$ and $W=A\cup B$.
For each ordered pair of distinct local vertices $p,q\in L$, introduce a Boolean variable $\alpha_{p,q}\in\{0,1\}$, where $\alpha_{p,q}=1$ means $p\to q$ ($q \in N^+(p)$).

\paragraph{Deterministic local constraints.}
The following constraints are imposed.
\begin{enumerate}[label=(R6.\arabic*)]
\item Orientation:
\[
\alpha_{p,q}+\alpha_{q,p}\le 1\qquad(p\ne q).
\]
\item The out-neighborhood of $a_1$ is fixed:
\[
\alpha_{a_1,v}=1\quad(v\in A_1\cup B),
\qquad
\alpha_{a_1,v}=0\quad(v\in X\cup R\cup C).
\]
\item Vertices of $A$ send no edge to $C$:
\[
\alpha_{a,c}=0\qquad(a\in A,\ c\in C).
\]
\item Every vertex of $A$ has at least two out-neighbors in $A$:
\[
\sum_{a'\in A\setminus\{a\}}\alpha_{a,a'}\ge 2\qquad(a\in A).
\]
\item Every witness has local outdegree at least seven:
\[
\sum_{v\in L\setminus\{w\}}\alpha_{w,v}\ge 7\qquad(w\in A\cup B).
\]
\item Vertices of $B$ do not send to $a_1$ or to $R$:
\[
\alpha_{b,a_1}=0,
\qquad
\alpha_{b,r}=0\qquad(b\in B,
\ r\in R).
\]
\item Vertices of $A_1$ do not send to $R$:
\[
\alpha_{u,r}=0\qquad(u\in A_1,
\ r\in R).
\]
\item Every vertex of $X$ is reached from $A_1\cup B$:
\[
\sum_{u\in A_1\cup B}\alpha_{u,x_i}\ge 1\qquad(1\le i\le x).
\]
\item Every vertex of $C$ is reached from $B$:
\[
\sum_{b\in B}\alpha_{b,c_j}\ge 1\qquad(1\le j\le M).
\]
\item The row-dependent counting lower bounds are imposed:
\[
\sum_{h\in H}\sum_{b\in B}\alpha_{h,b}\ge m_{HB}(x),
\qquad
\sum_{b\in B}\sum_{c\in C}\alpha_{b,c}\ge m_{BC}(x),
\]
where
\[
(m_{HB}(2),m_{HB}(3),m_{HB}(4))=(16,19,23),
\]
\[
(m_{BC}(2),m_{BC}(3),m_{BC}(4))=(19,16,14).
\]
\end{enumerate}

\paragraph{Local second-neighbor variables.}
For every witness $v\in(A\cup B)\setminus\{a_1\}$ and every $w\in L\setminus\{v\}$, introduce a Boolean variable
\[
\sigma_{v,w}\in\{0,1\}.
\]
The vertex $a_1$ is omitted because Lemma \ref{lem:r6-a1-exact} gives its second-neighborhood exactly.  The constraints are
\[
\sigma_{v,w}+\alpha_{v,w}\le 1,
\]
and, for every $u\in L\setminus\{v,w\}$,
\[
\sigma_{v,w}\ge \alpha_{v,u}+\alpha_{u,w}-\alpha_{v,w}-1.
\]
Equivalently, if $v\to u\to w$ and $v\not\to w$, then $\sigma_{v,w}=1$ is forced.  No converse implication is imposed: an unnecessary $\sigma_{v,w}$ is free to be zero.  Thus every genuine obstruction can assign $\sigma_{v,w}$ to be the exact local second-neighbor indicator, while the model remains a relaxation.

\paragraph{Outside-signature variables.}
For every nonempty subset $S\subseteq C$, introduce an integer variable
\[
n_S\in\{0,1,\ldots,7\}.
\]
It counts outside vertices $t\notin L$ whose exact predecessor set inside $C$ is $S$:
\[
\{c\in C:c\to t\}=S.
\]
The bound $n_S\le 7$ is harmless and valid: since $S$ is nonempty, each such outside vertex is an out-neighbor of every $c\in S$, and after trimming each $c$ has total outdegree seven.

The trimming equation for each $c\in C$ is
\begin{equation}\label{eq:r6-trim}
\sum_{v\in L\setminus\{c\}}\alpha_{c,v}+
\sum_{S\subseteq C:c\in S}n_S=7.
\end{equation}
For $b\in B$ and nonempty $S\subseteq C$, introduce a Boolean variable $\eta_{b,S}$ satisfying
\[
\eta_{b,S}=\bigvee_{c\in S}\alpha_{b,c}.
\]
Let $y_{b,S}$ be the linearized product $\eta_{b,S}n_S$; thus
\[
y_{b,S}=n_S\text{ if }b\text{ sends to at least one member of }S,
\qquad
 y_{b,S}=0\text{ otherwise}.
\]
\begin{remark}
Notice that for the auxiliary non-witness vertices whose outgoing edges are represented by signature variables, we trim their outgoing edges until their out-degree is exactly seven. But note that we could delete more if all we wanted was to preserve non-Seymourness of the witnesses. Deleting more outgoing edges from a non-witness would still only shrink second-neighborhoods of witnesses.

But for the finite model, deleting to exactly $7$ is better and more natural. If we deleted more than necessary, that would be a weaker and less precise model. A weaker model would allow for more artificial configurations and make infeasibility harder to prove. The equality of $7$ does not lose any genuine obstruction, because every outside non-witness vertex originally has at least $7$ outgoing edges.
\end{remark}
\paragraph{Non-Seymour inequalities.}
For every $a\in A\setminus\{a_1\}$, all possible second-neighbors are local, so the imposed inequality is
\begin{equation}\label{eq:r6-A-NS}
\sum_{w\in L\setminus\{a\}}\sigma_{a,w}
\le
\sum_{v\in L\setminus\{a\}}\alpha_{a,v}-1.
\end{equation}
For every $b\in B$, outside vertices reached through $C$ must also be counted:
\begin{equation}\label{eq:r6-B-NS}
\sum_{w\in L\setminus\{b\}}\sigma_{b,w}
+
\sum_{\varnothing\ne S\subseteq C}y_{b,S}
\le
\sum_{v\in L\setminus\{b\}}\alpha_{b,v}-1.
\end{equation}
The non-Seymour inequality for $a_1$ is exactly $x+M\le 7$, which is already built into the row list $R_6$.
\begin{remark}
    The model does not need to reconstruct the whole original graph. It needs only the information that can make a witness acquire a first or second out-neighbor. Vertices outside $L$ cannot be first out-neighbors of witnesses: all outgoing edges from $A\cup B$ are already in $L$. Such outside vertices can only be terminal second-neighbors reached from $C$, and their entire effect on the witness inequalities is captured by their predecessor subsets $S\subseteq C$.
\end{remark}

\begin{proposition}[Soundness of the $r=6$ model]\label{prop:r6-sound}
Every genuine $r=6$ obstruction gives a feasible assignment of the variables satisfying the finite model above for one of the five rows in $R_6$.
\end{proposition}
\begin{proof}
A genuine obstruction determines $x$ and $M$.  Lemmas~\ref{lem:r6-a1-exact} and \ref{lem:r6-counting} imply $(x,M)\in R_6$.  Now first trim the outgoing edges of the non-witness vertices in $C$, using Lemma~\ref{lem:trimming}, until each of them has total out-degree exactly seven.  This does not change any first-neighborhood of a witness and can only shrink witness second-neighborhoods.  Then assign the local arc variables $\alpha_{p,q}$ from the resulting local oriented graph on $L$.

Constraints (R6.1)--(R6.10) are exactly the definitions and counting inequalities above.  Count outside terminals by their exact predecessor set $S\subseteq C$; this gives the variables $n_S$ and the trimming equations \eqref{eq:r6-trim}.  The local second-neighbor constraints force every local directed two-step terminal that is not a first-neighbor to be counted.  Equations \eqref{eq:r6-A-NS} and \eqref{eq:r6-B-NS} are then exactly the non-Seymour inequalities for the witnesses after separating local second-neighbors from compressed outside terminals.  Hence the finite model is feasible.
\end{proof}


%\arp{After running computational checks for the model using OR-tools with the help of ChatGPT Pro, we get the following infeasibity result that proves existence of a Seymour vertex in $A\cup B$ for this branch. @Sandeep please check and edit accordingly}
The corresponding OR-Tools CP-SAT encodings were run to verify that all corresponding finite configurations are infeasible.
This implies the existence of a Seymour vertex in $A\cup B$ for this branch.

\begin{computationallemma}[Infeasibility of the $r=6$ residual rows]\label{cl:r6}
For every row $(x,M)\in R_6$, the exact outside-signature model described in Subsection \ref{subsec:r6-model} is infeasible.
\end{computationallemma}

\subsection{The branch $r=5$}
Now we have, $|N^+(a_1)\cap B|=5$. Let $P=N^+(a_1)\cap B$ and $B=P\cup\{q\}$. Thus, $N^+(a_1)=A_1\cup P$ and $|N^+(a_1)|=7$. In the $r=5$ branch we use the witness set $W=A\cup P$. The exceptional vertex $q$ is not a witness.

\begin{definition}
    Define $X=(N^+(A_1\cup P)\cap A)\setminus A_1$ and let $x=|X|$. Also let, $y=|N^+(A_1\cup P)\cap\{q\}|\in\{0,1\}$. 
\end{definition}

\begin{definition}
    Define $Z=N^+(P)\setminus(\{s\}\cup A\cup B)$ and let $z=|Z|$.
\end{definition}

Let $\varepsilon_s=1$ if some vertex of $P$ sends to $s$, and let $\varepsilon_s=0$ otherwise.

\begin{lemma}\label{lem:r5-a1}
In an $r=5$ obstruction,
\[
x+y+z+\varepsilon_s\le 6.
\]
\end{lemma}
\begin{proof}
The vertex $a_1$ has out-degree seven.  Its second-neighbors are exactly the vertices in $X$, possibly the vertex $q$ if $y=1$, the vertices in $Z$, and possibly the root $s$ if $\varepsilon_s=1$.  These vertices are distinct and none lies in $N^+(a_1)\cup\{a_1\}$.  Since $a_1$ is assumed to be non-Seymour as it belongs to the witness set, it has at most six second-neighbors. This finishes the proof.
\end{proof}



\begin{lemma}\label{lem:delete-q}
For the purpose of ruling out obstructions with witness set $A\cup P$, all outgoing edges of $q$ may be deleted.
\end{lemma}
\begin{proof}
Apply Lemma~\ref{lem:trimming} to the single non-witness vertex $q$. Deleting the outgoing edges of $q$ leaves all the first neighbors of witnesses unchanged and can only remove the second-neighbors of the witnesses.
\end{proof}

\begin{remark}
    In the \(r=5\) branch we test only the smaller set \(A\cup P\). Deleting all edges and not some edges of $q$ is just a choice to simplify the model. It weakens the model in a sound way and is sufficient for the final infeasibility check.
\end{remark}

\begin{definition} Define
    \[
H=A_1\cup X,
\qquad R=A\setminus(\{a_1\}\cup H),
\qquad h=|H|=x+2.
\]
\end{definition}

As before, Lemma~\ref{lem:X-at-least-2} gives $2\le x\le 4$.

\begin{lemma}[Counting in the $r=5$ branch]\label{lem:r5-counting}
In an $r=5$ obstruction, the only possible triples $(x,y,z)$ are
\begin{equation}\label{eq:R5rows}
R_5=\{(2,0,3),(2,0,4),(2,1,2),(2,1,3),(3,1,1),(3,1,2)\}.
\end{equation}
Moreover the value of $\varepsilon_s$ is forced as follows:
\[
\begin{array}{c|c}
(x,y,z)&\varepsilon_s\\ \hline
(2,0,3)&1\\
(2,0,4)&0\\
(2,1,2)&1\\
(2,1,3)&0\\
(3,1,1)&1\\
(3,1,2)&0
\end{array}
\]
\end{lemma}
\begin{proof}
No vertex of $A_1\cup P$ sends to $R$, and no vertex of $P$ sends to $a_1$.  Also vertices of $A_1$ do not send to $a_1$.

First lower-bound $e(H,P)$.  Vertices of $H$ have at least $7h$ outgoing edges into $A\cup B$.  At most
\[
\binom{h}{2}+x+x(4-x)
\]
of these edges end in $A$: the terms count edges inside $H$, edges from $X$ to $a_1$, and edges from $X$ to $R$.  The edges from $H$ to $q$ are at most $x$ when $y=0$, because then no vertex of $A_1$ sends to $q$, and at most $h=x+2$ when $y=1$.  Hence
\begin{equation}\label{eq:r5-eHP}
e(H,P)\ge 7h-\binom{h}{2}-x-x(4-x)-(x+2y).
\end{equation}

Now count outgoing edges of the five vertices in $P$.  Their out-neighbors lie in $H\cup P\cup Z\cup\{s,q\}$.  Therefore
\[
e(P,Z)+e(P,\{s\})+e(P,\{q\})
\ge 35-(5h-e(H,P))-10.
\]
Using \eqref{eq:r5-eHP}, this lower bound is
\[
\begin{array}{c|cc}
 x&y=0&y=1\\ \hline
 2&19&17\\
 3&16&14\\
 4&14&12
\end{array}
\]
On the other hand,
\[
e(P,Z)+e(P,\{s\})+e(P,\{q\})\le 5z+5\varepsilon_s+5y.
\]
Combining this with Lemma~\ref{lem:r5-a1}, namely
\[
z+\varepsilon_s\le 6-x-y,
\]
leaves only the possibilities
\[
\begin{array}{c|c|c}
x&y&z+\varepsilon_s\\ \hline
2&0&4\\
2&1&3\\
3&1&2.
\end{array}
\]
Since $\varepsilon_s\in\{0,1\}$, this gives exactly the row set and forced $\varepsilon_s$ values displayed above.
\end{proof}

\subsection*{The reduced $r=5$ model}
Fix a row $(x,y,z)\in R_5$ and the forced value of $\varepsilon_s$.  The explicit local set is $L=\{s\}\cup A\cup B\cup Z$, where
\[
A=\{a_1,u_1,u_2,x_1,\ldots,x_x,r_1,\ldots,r_{4-x}\},
\]
\[
P=\{p_1,\ldots,p_5\},
\qquad B=P\cup\{q\},
\qquad Z=\{z_1,\ldots,z_z\}.
\]
The witness set is $W=A\cup P$.
The model has Boolean variables $\alpha_{p,q}$ for all ordered pairs of distinct local vertices.  Its deterministic constraints are:
\begin{enumerate}[label=(R5.\arabic*)]
\item orientation: $\alpha_{p,q}+\alpha_{q,p}\le 1$ for $p\ne q$;
\item the root satisfies $N^+(s)=A$;
\item $a_1$ sends exactly to $A_1\cup P$;
\item all outgoing edges of $q$ are deleted;
\item vertices of $A$ do not send to $s$ or to $Z$;
\item every vertex of $A$ has at least two out-neighbors in $A$;
\item every witness in $A\cup P$ has local outdegree at least seven;
\item vertices of $P$ do not send to $a_1$ or to $R$;
\item vertices of $A_1$ do not send to $R$;
\item every vertex of $X$ is reached from $A_1\cup P$;
\item $y$ records whether $q$ is reached from $A_1\cup P$;
\item every vertex of $Z$ is reached from $P$;
\item the forced value of $\varepsilon_s$ records whether $s$ is reached from $P$.
\end{enumerate}
For every nonempty subset $S\subseteq Z$, introduce an integer signature variable $n_S\in\{0,\ldots,7\}$ counting outside vertices whose exact predecessor set in $Z$ is $S$.  After trimming the non-witness vertices of $Z$,
\begin{equation}\label{eq:r5-trim}
\sum_{v\in L\setminus\{z_i\}}\alpha_{z_i,v}
+
\sum_{S\subseteq Z:z_i\in S} n_S=7
\qquad(1\le i\le z).
\end{equation}
For each witness $v\in(A\cup P)\setminus\{a_1\}$, local second-neighbor variables and signature-hit variables are imposed exactly as in the $r=6$ model.  The resulting non-Seymour inequality is
\begin{equation}\label{eq:r5-NS}
\sum_{w\in L\setminus\{v\}}\sigma_{v,w}
+
\sum_{\varnothing\ne S\subseteq Z} y_{v,S}
\le
\sum_{u\in L\setminus\{v\}}\alpha_{v,u}-1.
\end{equation}
For $a_1$, the non-Seymour inequality is exactly $x+y+z+\varepsilon_s\le 6$, already built into the row list.

\begin{proposition}[Soundness of the $r=5$ model]\label{prop:r5-sound}
Every genuine $r=5$ obstruction gives a feasible assignment of the reduced finite model for one of the rows in $R_5$.
\end{proposition}
\begin{proof}
The obstruction determines $x,y,z$ and $\varepsilon_s$. Lemma \ref{lem:r5-counting} places $(x,y,z)$ in $R_5$ and forces $\varepsilon_s$.  Delete all outgoing edges of $q$, which is permitted by Lemma \ref{lem:delete-q}.  Then trim the outgoing edges of the non-witness vertices of $Z$ until each has total outdegree exactly seven.  Assign the local arc variables from this trimmed graph.  The deterministic constraints are precisely the definitions above.  The signature variables over $Z$ give \eqref{eq:r5-trim}, and the inequalities \eqref{eq:r5-NS} are exactly the non-Seymour inequalities for witnesses after local and compressed outside second-neighbors are separated.
\end{proof}

%\arp{please check :Now again after running the computations using OR-Tools with the help of ChatGPT Pro we found that all the residual rows were infeasible.}
The corresponding OR-Tools CP-SAT encodings were run to verify that all corresponding finite configurations are infeasible.
Thus we obtain Computational Lemma~\ref{cl:r5}.
\begin{computationallemma}[Infeasibility of the $r=5$ residual rows]\label{cl:r5}
For every row $(x,y,z)\in R_5$, the reduced root-aware outside-signature model described above is infeasible.
\end{computationallemma}




\begin{proof}[Proof of Theorem \ref{thm:main}]
Suppose, for contradiction, that no vertex of $A\cup B$ is Seymour.  Choose $a_1\in A$ with $|N^+(a_1)\cap A|=2$.

If $r=6$, the witness set is $A\cup B$.  By Proposition \ref{prop:r6-sound}, the assumed obstruction gives a feasible assignment of the $r=6$ finite model in one of the rows in $R_6$.  This contradicts Computational Lemma \ref{cl:r6}.

If $r=5$, write $P=N^+(a_1)\cap B$.  Since all vertices of $A\cup B$ are assumed non-Seymour, in particular all vertices of $A\cup P$ are non-Seymour.  By Proposition \ref{prop:r5-sound}, the assumed obstruction gives a feasible assignment of the reduced $r=5$ finite model in one of the rows in $R_5$.  This contradicts Computational Lemma \ref{cl:r5}.

Both possible values of $r$ lead to contradictions. Therefore at least one vertex in $A\cup B$ is Seymour.

\end{proof}


\section{Case  $|A|=7, |B|=6, |A_1|=1$}\label{sec:B6A11}

Let $A_1=\{u\}.$ Since $d^+(a_1)\ge 7$, $|A_1|=1$, and $|B|=6$, we have $N^+(a_1)=\{u\}\cup B$ and $d^+(a_1)=7$.

We assume that every vertex of $A\cup B$ is non-Seymour and derive an impossibility. The setting and terminology in this section is exactly same as Section~\ref{sec:B6A12}. We will use the exact same idea as before with slightly different constraints. We also have the exact same trimming lemma. In particular, when a non-witness vertex is needed only as an intermediate vertex for second-neighborhood counts, we may delete outgoing edges from that non-witness until its out-degree is exactly $7$. This operation never turns a non-Seymour witness into a Seymour witness.


The theorem we prove here is the following.

\begin{theorem}\label{thm:B6A11}
Let $D$ be a finite oriented graph with $\delta^+(D)=7$. Let $s$ be a vertex with $d^+(s)=7$, and let $A=N^+(s)$ and $B=N^{++}(s)$.
Assume $|A|=7$, $|B|=6$, and $\min_{a\in A}|N^+(a)\cap A|=1$.
Then some vertex in $A\cup B$ is a Seymour vertex.
\end{theorem}

Equivalently, in the notation $r=|N^+(a_1)\cap B|$, this entire case is the $r=6$ branch. The rest of the proof is by contradiction. We assume that every vertex of $A\cup B$ is non-Seymour and derive an impossibility.

\subsection*{Witnesses, obstructions, and trimming}

Recall that the vertices whose Seymour inequalities are imposed are called \emph{witnesses}. In this proof the witness set is $W=A\cup B$.
Also, recall that an \emph{obstruction} is a configuration satisfying all standing assumptions in which every vertex of $W$ is non-Seymour. The theorem says exactly that no obstruction exists.

The finite model below contains additional vertices outside $W$. These additional vertices are not tested for being Seymour vertices; they are only present because they may be first or second out-neighbors of vertices in $W$. Next we recall the trimming lemma below.

\begin{lemma}[Trimming lemma]\label{lem:recalltrimming}
Let $W$ be a witness set in an oriented graph $D$.  Let $D'$ be a subdigraph of $D$ obtained by deleting directed edges of $D$ whose tails all lie outside $W$.  If $v\in W$ is not a Seymour vertex in $D$, then $v$ is not a Seymour vertex in $D'$.
\end{lemma}

As before, this lemma is used only for non-witness vertices in the outside layer. After trimming, each such non-witness vertex may be assumed to have total outdegree exactly seven. This never turns a non-Seymour witness into a Seymour witness.

\begin{definition}
Define $X=\bigl(N^+(\{u\}\cup B)\cap A\bigr)\setminus\{u\}$ and
$x=|X|$. Also, define, $R=A\setminus(\{a_1,u\}\cup X).$
\end{definition}

Note that $a_1\to u$ and $a_1\to b$ for every $b\in B$. By the definition of $X$, no vertex of $\{u\}\cup B$ sends an edge to $R$. Hence, $N^+(\{u\}\cup B)\cap R=\emptyset.$ Also $|R|=5-x$.

Recall that $C=N^+(B)\setminus(A\cup B)$. Since $s\notin A\cup B$, this convention includes $s$ in $C$ exactly when some vertex of $B$ sends to $s$. Every vertex of $C$ is reached from $B$. Let $M=|C|$.
If $s\in C$, then $s$ is treated as an explicit non-witness auxiliary vertex. Its actual out-neighborhood is exactly $A$, and this assignment is allowed by the finite model.

\begin{lemma}\label{lem:xM}
In any obstruction, $1\le x\le 5$ and $x+M\le 6$.
Moreover, in the $r=6$ obstruction model, $N^{++}(a_1)=X\cup C$.
\end{lemma}

\begin{proof}
The upper bound $x\le 5$ follows from $X\subseteq A\setminus\{a_1,u\}$.

For the lower bound, the vertex $u$ has no out-neighbor in $\{a_1\}\cup R\cup C$. It cannot send to $a_1$ because $a_1\to u$; it cannot send to $R$ by the definition of $R$; and it cannot send to $C$ because $u\in A$ and every vertex of $A$ sends only into $A\cup B$. Thus all out-neighbors of $u$ lie in $X\cup B$. Since $d^+(u)\ge 7$ and $|B|=6$, we get $x\ge 1$.

Now consider second out-neighbors of $a_1$. The first out-neighbors of $a_1$ are exactly $u$ and the six vertices of $B$. By definition, the vertices of $X$ are precisely the vertices of $A\setminus\{u\}$ reached from $u$ or from $B$. Since neither $u$ nor any vertex of $B$ sends to $a_1$ or to $R$, the only second out-neighbors of $a_1$ in $A$ are the vertices of $X$. The vertices of $C$ are exactly the possible second out-neighbors of $a_1$ outside $A\cup B$ reached through $B$, with $s$ included exactly when it is reached through $B$. The vertices of $B$ and the vertex $u$ are first out-neighbors of $a_1$ and are excluded from $N^{++}(a_1)$. Hence $N^{++}(a_1)=X\cup C$.

Since $a_1$ is a witness vertex and thus assumed to be non-Seymour and $d^+(a_1)=7$, we have $|N^{++}(a_1)|\le 6$. Therefore $x+M\le 6$.
\end{proof}

Put $U=\{u\}\cup X$ and $|U|=x+1$.
For disjoint vertex sets $P,Q$, let $e(P,Q)$ denote the number of directed edges from $P$ to $Q$.

\begin{lemma}[Counting bounds]\label{lem:counting}
In any obstruction,
\[
e(U,B)\ge 7+\frac{x(x+1)}2
\]
and
\[
e(B,C)\ge 28-6x+\frac{x(x+1)}2.
\]
Consequently:
\[
\begin{array}{c|ccccc}
x&1&2&3&4&5\\ \hline
e(U,B)\ge&8&10&13&17&22\\
e(B,C)\ge&23&19&16&14&13
\end{array}
\]
\end{lemma}

\begin{proof}
Let $h=|U|=x+1$. Since $U\subseteq A$, every vertex of $U$ sends all its out-neighbors into $A\cup B$, and every vertex of $U$ has outdegree at least seven. Thus
\begin{equation}
e(U,A)+e(U,B)\ge 7h.\tag{1}
\end{equation}
We upper-bound $e(U,A)$. Inside $U$ there are at most $\binom h2$ oriented edges. The vertex $u$ sends no edge to $a_1$ or to $R$. Each of the $x$ vertices of $X$ may send one edge to $a_1$, and may send edges to at most $5-x$ vertices of $R$. Therefore
\[
e(U,A)\le \binom h2+x+x(5-x).
\]
Together with (1), this gives
\[
e(U,B)\ge 7h-\left(\binom h2+x+x(5-x)\right)
=7+\frac{x(x+1)}2.
\]

Next count outgoing edges from the six vertices of $B$. No vertex of $B$ sends to $a_1$, because $a_1\to B$, and no vertex of $B$ sends to $R$ by the definition of $R$. Hence all out-neighbors of vertices of $B$ lie in $U\cup B\cup C$. Therefore
\[
42\le e(B,U)+e(B,B)+e(B,C).
\]
Between $B$ and $U$ there are $6h$ unordered cross-pairs, and the graph is oriented, so
\[
e(B,U)\le 6h-e(U,B).
\]
Also $e(B,B)\le \binom 62=15$. Hence
\[
e(B,C)\ge 42-(6h-e(U,B))-15.
\]
Substituting the lower bound on $e(U,B)$ gives
\[
e(B,C)\ge 28-6x+\frac{x(x+1)}2.
\]
The displayed table follows by substituting $x=1,2,3,4,5$.
\end{proof}

Since each vertex of $C$ receives at most one edge from each of the six vertices of $B$, we have,
\begin{equation}
e(B,C)\le 6M.\tag{2}
\end{equation}
Combining Lemmas~\ref{lem:xM} and \ref{lem:counting} with (2) gives the residual rows.

\begin{proposition}[Residual rows]\label{prop:rows}
Every obstruction in the present case satisfies
\[
(x,M)\in R_1:=\{(1,4),(1,5),(2,4),(3,3)\}.
\]
\end{proposition}

\begin{proof}
By Lemma \ref{lem:xM}, $1\le x\le 5$ and $M\le 6-x$. By Lemma \ref{lem:counting} and (5), we also need
\[
6M\ge 28-6x+\frac{x(x+1)}2.
\]
Checking $x=1,2,3,4,5$ gives:
\[
\begin{array}{c|c|c|c}
x&M\le 6-x&e(B,C)\ge&\text{surviving }M\\ \hline
1&5&23&4,5\\
2&4&19&4\\
3&3&16&3\\
4&2&14&\text{none}\\
5&1&13&\text{none}
\end{array}
\]
This is exactly the stated row set.
\end{proof}

\subsection*{The exact finite-signature obstruction model}

This section gives the exact finite model used in the computation. Fix one residual row $(x,M)\in R_1$. The explicit local vertex set is $L=A\cup B\cup C$, where,
\[
A=\{a_1,u\}\cup X\cup R,
\qquad |X|=x,
\qquad |R|=5-x,
\qquad |B|=6,
\qquad |C|=M.
\]
The witness set is $W=A\cup B$. The vertices of $C$ are non-witness vertices. For every ordered pair of distinct vertices $p,q\in L$, introduce a Boolean variable
\[
a_{p,q}\in\{0,1\},
\]
with $a_{p,q}=1$ meaning $p\to q$. Orientation is imposed by
\begin{equation}
a_{p,q}+a_{q,p}\le 1\qquad (p\ne q).\tag{3}
\end{equation}

The deterministic constraints are as follows.
\begin{enumerate}[label=(F\arabic*),leftmargin=*]
\item The special vertex $a_1$ sends exactly to $u$ and to all six vertices of $B$:
\[
\begin{aligned}
a_{a_1,u}&=1,\\
a_{a_1,b}&=1 &&(b\in B),\\
a_{a_1,v}&=0 &&(v\in X\cup R\cup C).
\end{aligned}
\]
\item No vertex of $A$ sends to $C$:
\[
a_{a,c}=0\qquad (a\in A,
\ c\in C).
\]
\item No vertex of $\{u\}\cup B$ sends to $a_1$ or to $R$:
\[
a_{v,a_1}=0\qquad (v\in\{u\}\cup B),
\]
\[
a_{v,r}=0\qquad (v\in\{u\}\cup B,
\ r\in R).
\]
\item Each vertex of $X$ is reached from $\{u\}\cup B$:
\[
\sum_{v\in\{u\}\cup B}a_{v,x_i}\ge 1\qquad (x_i\in X).
\]
\item Each vertex of $C$ is reached from $B$:
\[
\sum_{b\in B}a_{b,c_j}\ge 1\qquad (c_j\in C).
\]
\item Every vertex of $A$ has at least one out-neighbor in $A$:
\[
\sum_{t\in A\setminus\{a\}}a_{a,t}\ge 1\qquad (a\in A).
\]
This records $\min_{a\in A}|N^+(a)\cap A|=1$ together with the fact that $a_1$ has exactly one such neighbor.
\item Every witness has local outdegree at least seven:
\[
\sum_{t\in L\setminus\{w\}}a_{w,t}\ge 7\qquad (w\in A\cup B).
\]
For witnesses this local outdegree is the actual outdegree, because all outgoing edges from $A\cup B$ lie inside $L$.
\item The handwritten aggregate lower bounds are imposed:
\[
E_{UB}:=\sum_{u'\in U}\sum_{b\in B}a_{u',b}
\ge 7+\frac{x(x+1)}2,
\]
\[
E_{BC}:=\sum_{b\in B}\sum_{c\in C}a_{b,c}
\ge 28-6x+\frac{x(x+1)}2.
\]
\end{enumerate}

Outside vertices matter only as terminal second-neighbors reached from $C$. They are compressed by exact predecessor signatures. For every nonempty subset $S\subseteq C$, introduce an integer variable
\[
n_S\in\{0,1,\ldots,7\},
\]
where $n_S$ is the number of outside vertices $t\notin L$ whose set of predecessors in $C$ is exactly $S$:
\[
\{c\in C:c\to t\}=S.
\]
After trimming every vertex of $C$ to total outdegree exactly seven, the model imposes
\begin{equation}
d_L^+(c)+\sum_{S\ni c}n_S=7\qquad(c\in C),\tag{4}
\end{equation}
where
\[
d_L^+(c)=\sum_{t\in L\setminus\{c\}}a_{c,t}.
\]

For each witness $v\in W$ and each $w\in L\setminus\{v\}$, introduce a Boolean variable
\[
\sigma_{v,w}\in\{0,1\},
\]
which records whether $w$ is forced to be a local second out-neighbor of $v$. First out-neighbors are excluded by
\begin{equation}
\sigma_{v,w}+a_{v,w}\le 1.\tag{5}
\end{equation}
For every possible intermediate vertex $y\in L\setminus\{v,w\}$, the implication
\[
(v\to y)\wedge(y\to w)\wedge(v\not\to w)\Longrightarrow \sigma_{v,w}=1
\]
is encoded linearly as
\begin{equation}
\sigma_{v,w}\ge a_{v,y}+a_{y,w}-a_{v,w}-1.\tag{6}
\end{equation}
Equivalently, the implementation uses the Boolean clause
\[
\neg a_{v,y}\vee \neg a_{y,w}\vee a_{v,w}\vee \sigma_{v,w}.
\]

For $a\in A$, there is no first step from $a$ into $C$, so all possible second-neighbors of $a$ are already local. The non-Seymour inequality is
\begin{equation}
\sum_{w\in L\setminus\{a\}}\sigma_{a,w}
\le d_L^+(a)-1.
\tag{7}
\end{equation}
For $b\in B$, outside signature variables may contribute. For every nonempty $S\subseteq C$, introduce a Boolean hit variable $h_{b,S}$ with
\begin{equation}
h_{b,S}=1\quad\Longleftrightarrow\quad \sum_{c\in S}a_{b,c}\ge 1.
\tag{8}
\end{equation}
The implementation enforces (8) by $a_{b,c}\le h_{b,S}$ for all $c\in S$ and $h_{b,S}\le \sum_{c\in S}a_{b,c}$. It then introduces an integer product variable
\[
p_{b,S}=h_{b,S}n_S
\]
using the standard bounds
\[
0\le p_{b,S}\le 7h_{b,S},\qquad p_{b,S}\le n_S,
\qquad p_{b,S}\ge n_S-7(1-h_{b,S}).
\]
The non-Seymour inequality for $b\in B$ is
\begin{equation}
\sum_{w\in L\setminus\{b\}}\sigma_{b,w}
+\sum_{\emptyset\ne S\subseteq C}p_{b,S}
\le d_L^+(b)-1.
\tag{9}
\end{equation}

%\begin{remark}[Why this model is exact enough]
%The model does not need to reconstruct the whole original graph. It needs only the information that can make a witness acquire a first or second out-neighbor. Vertices outside $L$ cannot be first out-neighbors of witnesses: all outgoing edges from $A\cup B$ are already in $L$. Such outside vertices can only be terminal second-neighbors reached from $C$, and their entire effect on the witness inequalities is captured by their predecessor subsets $S\subseteq C$.
%\end{remark}

\begin{lemma}[Soundness of the finite model]\label{lem:sound}
If an obstruction exists in the original graph for a residual row $(x,M)$, then the corresponding finite exact-signature model for $(x,M)$ is feasible.
\end{lemma}

\begin{proof}
Start with an obstruction. Label $a_1,u,X,R,B,C$ as above. If $s$ is absorbed into $C$, regard it as one of the non-witness vertices of $C$; its actual outgoing edges are exactly the seven edges from $s$ to $A$.

Apply Lemma~\ref{lem:recalltrimming} to delete outgoing edges from non-witness vertices until each vertex of $C$ has total outdegree exactly seven. No edge leaving a witness is deleted, and every witness remains non-Seymour. Now record the remaining directed edges inside $L$ as Boolean variables $a_{p,q}$. Conditions (F1)--(F7) follow directly from the definitions of $A,B,X,R,C$ and from $\delta^+(D)=7$; condition (F8) follows from Lemma~\ref{lem:counting}.

For each outside out-neighbor $t\notin L$ of a vertex of $C$, record the nonempty predecessor set $S_t=\{c\in C:c\to t\}$ and add $t$ to the corresponding count $n_{S_t}$. This gives exactly equation (4). The variables $\sigma_{v,w}$ record all forced local second-neighbors, and the hit/product variables record exactly which compressed outside terminals are seen from each $b\in B$. Thus, inequalities (7) and (9) are precisely the non-Seymour inequalities for the witnesses after trimming. Therefore, the finite model is feasible.
\end{proof}
\begin{computationallemma}[Finite infeasibility certificate]\label{lem:comp}
The finite exact-signature model is infeasible for every residual row in $R_1$.
\end{computationallemma}


\begin{proof}[Proof of Theorem~\ref{thm:B6A11}]
Assume for contradiction that every vertex of $A\cup B$ is non-Seymour. Choose $a_1\in A$ with $|N^+(a_1)\cap A|=1$, write $N^+(a_1)\cap A=\{u\}$, and define $X,R,C,x,M$ as above. Proposition \ref{prop:rows} shows that
\[
(x,M)\in\{(1,4),(1,5),(2,4),(3,3)\}.
\]
By Lemma~\ref{lem:sound}, the assumed obstruction gives a feasible finite exact-signature model for this row. But Lemma~\ref{lem:comp} says that every such finite model is infeasible. This contradiction proves that not all vertices of $A\cup B$ are non-Seymour. Therefore, at least one vertex in $A\cup B$ is a Seymour vertex.
\end{proof}


\appendix

\section{Missing proofs of preliminaries}\label{sec:Missing proofs in preliminaries}

\begin{lemma}\label{app:lem:general bound for A_1}

\[
|A_{1}|\leq \left\lceil \frac{\delta}{2}\right\rceil-1.
\]
\end{lemma}

\begin{proof}
By the choice of \(a_{1}\), every vertex \(a\in A\) has at least \(|A_{1}|\) out-neighbours inside \(A\). Therefore the number of arcs in the induced oriented graph \(D[A]\) is at least $|A|\cdot |A_{1}|.$
On the other hand, since \(D\) is a simple oriented graph, between every pair of vertices of \(A\) there is at most one arc. Hence the number of arcs in \(D[A]\) is at most $\binom{|A|}{2}.$
Therefore
\[
|A|\cdot |A_{1}|\leq \binom{|A|}{2}.
\]
Since \(|A|=\delta\), this gives
\[
\delta |A_{1}|\leq \frac{\delta(\delta-1)}{2}.
\]
As \(\delta>0\), and  \(|A_{1}|\) is an integer, we have our lemma.

\end{proof}


\begin{lemma}~\label{app:lem:perfectdegree}
Suppose that \(|A|\) is odd and
\[
|A_{1}|=\frac{|A|-1}{2}.
\]
Then D[A] is a regular tournament and $|A_2|=|A_{1}|.$
\end{lemma}

\begin{proof}
Let $|A|=2k+1$, thus \(|A_{1}|=k\). By counting the number of arcs in \(D[A]\) ($a_1$ has minimum-outdegree) it is easy to observe that D[A] must be a regular tournament where each vertex of $A$ has exactly \(k\) out-neighbours inside \(A\).

Now consider a vertex $x\in A\setminus(\{a_{1}\}\cup A_{1}\cup A_{2}).$
Since \(x\notin A_{1}\), we do not have \(a_{1}x\in E(D)\). Since \(D[A]\) is a tournament, this implies $xa_{1}\in E(D).$

Also, since \(x\notin A_{2}\), no vertex of \(A_{1}\) sends an arc to \(x\). Hence for every \(y\in A_{1}\), because \(D[A]\) is a tournament, we must have $xy\in E(D).$

Thus \(x\) sends arcs to \(a_{1}\) and to every vertex of \(A_{1}\). Therefore \(x\) has at least $1+|A_{1}|>k$ out-neighbours inside \(A\), contradicting the fact that every vertex of \(A\) has exactly \(k\) out-neighbours inside \(A\).

Hence no such \(x\) exists. Therefore 
\[
A=\{a_{1}\}\cup A_{1}\cup A_{2}
\]

Hence the lemma follows.

\end{proof}

\begin{proposition}~\label{app: prop:easyseymour}

\begin{enumerate}
    \item If $|B|\geq \delta$, then $s$ is a Seymour vertex.

\item If $|B|\leq \left\lceil \frac{\delta}{2}\right\rceil$
and \(\delta\geq 2\), then $a_1$ is a Seymour vertex.
\end{enumerate}
\end{proposition}

\begin{proof}
Note that the first statement is trivial. Now suppose that $|B|\leq \left\lceil \frac{\delta}{2}\right\rceil.$ First assume that \(\delta=2k\) is even. Then $|B|\leq k.$
By Lemma 1,
\[
|A_{1}|\leq \left\lceil \frac{2k}{2}\right\rceil-1=k-1.
\]
Since all out-neighbours of \(a_{1}\in A\) lie in \(A\cup B\), we have, $d_{+}(a_{1})\leq |A_{1}|+|B|.$
Thus $d_{+}(a_{1})\leq (k-1)+k=2k-1<2k=\delta$,
contradicting the definition of \(\delta\) as the minimum outdegree. Therefore the even case cannot occur.

Now assume that $\delta=2k+1$ is odd. Since \(\delta\geq 2\), we have \(k\geq 1\). The hypothesis gives $|B|\leq k+1.$
By Lemma 1, $|A_{1}|\leq \left\lceil \frac{2k+1}{2}\right\rceil-1=k.$

Again, $d_{+}(a_{1})\leq |A_{1}|+|B|\leq k+(k+1)=2k+1=\delta$.
But \(d_{+}(a_{1})\geq \delta\), because \(\delta\) is the minimum outdegree. Hence equality must hold throughout. Therefore
\[
|A_{1}|=k,
\qquad
|B|=k+1,
\qquad
N^{+}(a_{1})\cap B=B,
\qquad
d_{+}(a_{1})=2k+1.
\]
In particular, \(a_{1}\) dominates \(B\).

Moreover, since \(a_{1}\) was chosen with the minimum possible number of out-neighbours in \(A\), every vertex \(a\in A\) has at least \(|A_{1}|=k\) out-neighbours inside \(A\). Since \(|B|=k+1\) and every vertex has outdegree at least \(2k+1\), each vertex of \(A\) must have exactly \(k\) out-neighbours in \(A\) and must dominate all of \(B\). Hence \(A\) dominates \(B\).

Because \(A\) dominates \(B\), there are no arcs from \(B\) to \(A\), since the graph is oriented and has no directed \(2\)-cycles. Thus every out-neighbour of a vertex of \(B\) lies in \(B\cup C\).

The total number of arcs leaving \(B\) is at least $\delta |B|.$

The number of arcs with both ends in \(B\) is at most $\binom{|B|}{2}.$
Therefore the number of arcs from \(B\) to \(C\) is at least
\[
\delta |B|-\binom{|B|}{2}.
\]
Each vertex of \(C\) can receive arcs from at most \(|B|\) vertices of \(B\). Hence the number of vertices of \(C\) dominated by \(B\) is at least
\[
\left\lceil
\frac{\delta |B|-\binom{|B|}{2}}{|B|}
\right\rceil.
\]
Using \(\delta=2k+1\) and \(|B|=k+1\), this becomes
\[
\left\lceil
\frac{(2k+1)(k+1)-\binom{k+1}{2}}{k+1}
\right\rceil
=
\left\lceil
2k+1-\frac{k}{2}
\right\rceil
=
\left\lceil \frac{3k}{2}\right\rceil+1.
\]
Since \(a_{1}\) dominates \(B\), every vertex of \(C\) dominated by \(B\) is at out-distance exactly \(2\) from \(a_{1}\). Thus
\[
|N^{++}(a_{1}) \cap C|\geq \left\lceil \frac{3k}{2}\right\rceil+1.
\]

Also, \(|A|=\delta=2k+1\) and \(|A_{1}|=k=(|A|-1)/2\). Thus, by Lemma~\ref{app:lem:perfectdegree}, $|N^{++}(a_{1}) \cap A|=|A_{2}|=k$.
Since \(A\cap C=\emptyset\) and \(k\geq 1\), by adding the size of the second neighborhood of $a_1$ inside $A$ and $C$, it is easy to see that $a_1$ is a Seymour vertex.

\end{proof}






\section*{Acknowledgements}
The author Arpan Sadhukhan gratefully acknowledges support from the Anusandhan National Research Foundation (ANRF) through the ANRF-National Post Doctoral Fellowship (N-PDF), File No. PDF/2025/001644.  The authors acknowledge the assistence from ChatGPT 5.5 Pro for the development of the CP-SAT models used in the computational part; the resulting OR-Tools CP-SAT encodings were then run and independently audited by the authors to verify that all corresponding finite configurations are infeasible. The files can be found at \url{https://github.com/rbsandeep/Seymour-Vertex-delta7}. The authors also acknowledge the use of ChatGPT 5.5 Plus and Pro models for mathematical checks for any possible errors in parts of the paper.



\section*{Data availability}

No data were used for the research described in this article.

\bibliographystyle{model1-num-names}
\bibliography{refs}

\end{document}
